\documentclass[11pt]{article}
\usepackage{cornell-notes}
\usepackage{mathtools}
\usepackage{needspace}
\usepackage{ragged2e}
\tcbuselibrary{most}
\geometry{letterpaper,margin=0.72in,headheight=15pt}
\setCornellCollection{Cloud Architecture Cornell Notes}
\setCornellUnitType{Chapter}
\setCornellUnitNumber{16}
\setCornellUnitTitle{Cloud Delivery Model Considerations}

\definecolor{Navy}{HTML}{17324D}
\definecolor{Teal}{HTML}{147D78}
\definecolor{Sky}{HTML}{E8F2F7}
\definecolor{CueGray}{HTML}{EEF2F4}
\definecolor{Warm}{HTML}{FFF3D6}
\definecolor{Risk}{HTML}{953B35}
\definecolor{Green}{HTML}{276749}
\definecolor{RuleGray}{HTML}{AAB8C2}
\definecolor{Text}{HTML}{25313A}

\hypersetup{colorlinks=true,linkcolor=Navy,urlcolor=Teal,citecolor=Navy,pdfborder={0 0 0}}
\color{Text}
\setlength{\parindent}{0pt}
\setlength{\parskip}{4pt}
\setlength{\emergencystretch}{3em}
\setlength{\LTpre}{3pt}
\setlength{\LTpost}{5pt}
\renewcommand{\arraystretch}{1.16}
\setlist{nosep,leftmargin=1.45em}
\hfuzz=0.2pt
\hbadness=2500

\titleformat{\section}{\Large\bfseries\color{Navy}\RaggedRight}{}{0pt}{}[\color{Teal}\titlerule]
\titleformat{\subsection}{\large\bfseries\color{Teal}\RaggedRight}{}{0pt}{}
\titlespacing*{\section}{0pt}{1.35ex plus .4ex}{.7ex}
\titlespacing*{\subsection}{0pt}{1.1ex plus .3ex}{.45ex}

\pagestyle{fancy}
\fancyhf{}
\fancyhead[L]{\small\color{Navy}\textbf{Cloud Architecture Cornell Notes}}
\fancyhead[R]{\small\color{Teal}\ChapterMark}
\fancyfoot[L]{\footnotesize\color{Navy}Concepts $\bullet$ Mechanisms $\bullet$ Evidence $\bullet$ Decisions}
\fancyfoot[R]{\footnotesize\color{Navy}\thepage\ of \pageref*{LastPage}}
\renewcommand{\headrulewidth}{0.5pt}
\renewcommand{\footrulewidth}{0.35pt}

\newtcolorbox{orientationbox}{enhanced,breakable,colback=Sky,colframe=Teal,boxrule=.75pt,arc=2mm,title={Chapter Orientation},fonttitle=\bfseries,left=2mm,right=2mm,top=1.5mm,bottom=1.5mm}
\newtcolorbox{sectionsummarybox}{enhanced,breakable,colback=CueGray,colframe=RuleGray,boxrule=.55pt,arc=1.5mm,title={Topic Summary},fonttitle=\bfseries\color{Navy},left=2mm,right=2mm,top=1mm,bottom=1mm}
\newtcolorbox{practicebox}{enhanced,breakable,colback=white,colframe=Green,boxrule=.75pt,arc=2mm,title={Practitioner Lens},fonttitle=\bfseries}
\newtcolorbox{scopenotebox}{enhanced,breakable,colback=Warm,colframe=Risk,boxrule=.7pt,arc=2mm,title={Terminology and Scope Note},fonttitle=\bfseries}
\newtcolorbox{synthesisbox}{enhanced,breakable,colback=Navy!5,colframe=Navy,boxrule=.85pt,arc=2mm,title={Chapter Synthesis},fonttitle=\bfseries}
\newtcolorbox{takeawaybox}{enhanced,breakable,colback=Teal!8,colframe=Teal,boxrule=.9pt,arc=2mm,title={One-Sentence Takeaway},fonttitle=\bfseries}
\newcommand{\CornellHeader}{\rowcolor{Navy}\color{white}\bfseries Cue / Active Recall & \color{white}\bfseries Notes / Reasoning \\\endhead}
\newcommand{\CornellRow}[2]{\cellcolor{CueGray}\textbf{#1} & #2 \\\addlinespace[3pt]}

\newcommand{\ChapterMark}{Chapter 16}
\title{Chapter 16: Cloud Delivery Model Considerations\\[2mm]\large Cornell Notes}
\author{}
\date{}
\hypersetup{pdftitle={Chapter 16: Cloud Delivery Model Considerations - Cornell Notes},pdfauthor={},pdfsubject={Cloud computing and architecture study notes},pdfkeywords={cloud, architecture, security, operations, Cornell notes}}
\begin{document}
\maketitle
\thispagestyle{fancy}
\vspace{-1.5em}
\begin{orientationbox}
\textbf{Central problem.} Compare provider and consumer responsibilities, investments, controls, and operating concerns for infrastructure, platform, and software service delivery.\par\smallskip
\textbf{Prerequisites.} Cloud delivery models, data centers, virtualization, platforms, applications, monitoring, security, service management, pricing, and SLAs.\par\smallskip
\textbf{Connection.} This chapter applies the earlier models and architectures to practical service construction and consumption before cost and quality are quantified.
\end{orientationbox}
\section{Learning Objectives}
\begin{enumerate}
\item Define the principal concepts and decision boundaries in Cloud Delivery Model Considerations.
\item Explain how 16.1 Provider perspective - building IaaS contributes to the chapter's central problem.
\item Distinguish Data centers and resource pools from Scalability and reliability using responsibility, behavior, and evidence.
\item Trace components, data, control, identity, capacity, and failure across 16.1 Provider perspective - equipping PaaS.
\item Analyze benefits together with assumptions, limitations, security risks, and operational tradeoffs.
\item Evaluate provider and consumer responsibilities at every relevant trust or administrative boundary.
\item Select verification evidence for architecture, configuration, monitoring, cost, and service-quality claims.
\item Apply the chapter's reasoning to a realistic design, review, migration, or incident scenario.
\end{enumerate}
\section{Cornell Cue-and-Notes}
\Needspace{9\baselineskip}
\subsection{16.1 Provider perspective - building IaaS}
\begin{longtable}{@{}C N@{}}
\CornellHeader
\CornellRow{What should be remembered about Data centers and resource pools?}{The provider invests in physical sites and organizes standardized compute, storage, and network resources into managed pools exposed through virtual infrastructure. \textbf{Review lens.} Plan capacity, regions, failure domains, lifecycle, utilization, reservations, spares, physical security, and environmental resilience.}
\CornellRow{What should be remembered about Scalability and reliability?}{Replication, clustering, load distribution, elastic allocation, and failover enable variable capacity and continuity. \textbf{Review lens.} Publish realistic limits and test saturation, correlated failure, maintenance, recovery, and performance after loss.}
\CornellRow{What should be remembered about Monitoring and management?}{Usage, health, audit, SLA, resource, and billing systems operate the service and create consumer-facing evidence. \textbf{Review lens.} Define telemetry coverage, tenant attribution, retention, reporting, alert ownership, and protection of the control plane.}
\CornellRow{What should be remembered about IaaS security?}{The provider protects facilities, physical resources, virtualization, core networks, service interfaces, and provider identities while exposing configuration controls to consumers. \textbf{Review lens.} Make the shared-responsibility boundary explicit and supply evidence that consumers can use to govern their layers.}
\end{longtable}
\begin{sectionsummarybox}
An IaaS provider assembles facilities, compute, virtualization, storage, networking, automation, monitoring, security, and support into a consumable infrastructure service.
\end{sectionsummarybox}
\Needspace{9\baselineskip}
\subsection{16.1 Provider perspective - equipping PaaS}
\begin{longtable}{@{}C N@{}}
\CornellHeader
\CornellRow{What should be remembered about Managed platform?}{The platform standardizes supported languages, runtimes, frameworks, services, deployment units, configuration, scaling, and observability. \textbf{Review lens.} Control runtime versions, patching, tenant isolation, build provenance, secrets, service bindings, quotas, and deprecation.}
\CornellRow{What should be remembered about Platform scalability and reliability?}{The provider scales runtime instances and supporting platform services while hiding much infrastructure management from consumers. \textbf{Review lens.} Expose platform limits and dependency behavior; test data, messaging, identity, deployment, and control-plane failure.}
\CornellRow{What should be remembered about PaaS monitoring and security?}{Platform telemetry and controls must cover application deployment, runtime, dependencies, service identities, data, and consumer actions. \textbf{Review lens.} Offer actionable logs and metrics, secure deployment pipelines, isolate tenants, and define which application vulnerabilities remain consumer-owned.}
\end{longtable}
\begin{sectionsummarybox}
A PaaS provider adds managed runtimes, middleware, development and deployment services, data services, and application lifecycle capabilities above infrastructure.
\end{sectionsummarybox}
\Needspace{9\baselineskip}
\subsection{16.1 Provider perspective - optimizing SaaS}
\begin{longtable}{@{}C N@{}}
\CornellHeader
\CornellRow{What should be remembered about SaaS operating model?}{The provider manages application code, platform, infrastructure, releases, data services, monitoring, support, and continuity while exposing tenant configuration and use. \textbf{Review lens.} Design tenant isolation, release compatibility, configuration governance, data lifecycle, export, identity federation, support access, and incident communication.}
\CornellRow{What should be remembered about SaaS security?}{Provider control increases, but consumer data, users, entitlements, integrations, and business configuration still determine significant risk. \textbf{Review lens.} Supply granular authorization, audit evidence, encryption, secure defaults, retention, deletion, backup, recovery, and transparent subprocessors.}
\end{longtable}
\begin{sectionsummarybox}
A SaaS provider operates the full application stack and must deliver secure, scalable, configurable, supportable multitenant business functionality.
\end{sectionsummarybox}
\Needspace{9\baselineskip}
\subsection{16.2 Consumer perspective - working with IaaS}
\begin{longtable}{@{}C N@{}}
\CornellHeader
\CornellRow{What should be remembered about IaaS provisioning decisions?}{The consumer selects regions, images, instance sizes, networks, storage, access controls, availability patterns, and lifecycle automation. \textbf{Review lens.} Use approved templates, policy enforcement, least privilege, patching, inventory, monitoring, backup, cost controls, and deprovisioning.}
\end{longtable}
\begin{sectionsummarybox}
An IaaS consumer gains low-level flexibility but retains broad responsibility for virtual networks, systems, middleware, applications, identities, data, and operations.
\end{sectionsummarybox}
\Needspace{9\baselineskip}
\subsection{16.2 Consumer perspective - working with PaaS}
\begin{longtable}{@{}C N@{}}
\CornellHeader
\CornellRow{What should be remembered about PaaS provisioning decisions?}{The consumer chooses runtime and service plans, application topology, scaling, bindings, configuration, identity, data services, and deployment process. \textbf{Review lens.} Assess portability, supported versions, quotas, observability, build security, secret handling, data protection, resilience, and platform exit.}
\end{longtable}
\begin{sectionsummarybox}
A PaaS consumer deploys applications to a managed runtime and trades infrastructure control for faster delivery and platform constraints.
\end{sectionsummarybox}
\Needspace{9\baselineskip}
\subsection{16.2 Consumer perspective - working with SaaS}
\begin{longtable}{@{}C N@{}}
\CornellHeader
\CornellRow{What should be remembered about SaaS consumption decisions?}{The consumer maps business processes and data to a provider-operated application and configures roles, policies, integrations, and retention. \textbf{Review lens.} Review functional fit, identity, authorization, data location, audit export, availability, support, privacy, customization, lock-in, continuity, and termination.}
\end{longtable}
\begin{sectionsummarybox}
A SaaS consumer governs tenant configuration, users, access, data, integration, workflow, quality, vendor risk, and exit rather than the application implementation.
\end{sectionsummarybox}
\Needspace{9\baselineskip}
\subsection{16.3 Case application}
\begin{longtable}{@{}C N@{}}
\CornellHeader
\CornellRow{What should be remembered about Provider-consumer responsibility matrix?}{For every stack layer and lifecycle activity, identify who configures, operates, secures, monitors, backs up, restores, supports, and supplies evidence. \textbf{Review lens.} Reconcile the matrix with contracts, architecture, platform behavior, internal staffing, and incident procedures.}
\end{longtable}
\begin{sectionsummarybox}
A delivery-model decision should compare responsibility, control, speed, skill, cost, portability, and quality rather than assuming one model is universally superior.
\end{sectionsummarybox}
\begin{practicebox}
For each model, mechanism, architecture, metric, or case decision, record the objective, owner, trust boundary, inputs, dependencies, expected behavior, failure consequence, measurement point, response, and retained evidence. A diagram or product label is a hypothesis until runtime and recovery tests support it.
\end{practicebox}
\section{Chapter Recap Matrix}
\begingroup\small\setlength{\tabcolsep}{4pt}
\begin{longtable}{@{}>{\RaggedRight\arraybackslash}p{0.235\textwidth}>{\RaggedRight\arraybackslash}p{0.37\textwidth}>{\RaggedRight\arraybackslash}p{0.325\textwidth}@{}}
\toprule\rowcolor{Navy}\color{white}\textbf{Element} & \color{white}\textbf{Purpose or Behavior} & \color{white}\textbf{Decision and Evidence}\\\midrule\endfirsthead
\toprule\rowcolor{Navy}\color{white}\textbf{Element} & \color{white}\textbf{Purpose or Behavior} & \color{white}\textbf{Decision and Evidence}\\\midrule\endhead
Data centers and resource pools & The provider invests in physical sites and organizes standardized compute, storage, and network resources into managed pools exposed through virtual infrastructure. & Plan capacity, regions, failure domains, lifecycle, utilization, reservations, spares, physical security, and environmental resilience. \\\midrule
Scalability and reliability & Replication, clustering, load distribution, elastic allocation, and failover enable variable capacity and continuity. & Publish realistic limits and test saturation, correlated failure, maintenance, recovery, and performance after loss. \\\midrule
Monitoring and management & Usage, health, audit, SLA, resource, and billing systems operate the service and create consumer-facing evidence. & Define telemetry coverage, tenant attribution, retention, reporting, alert ownership, and protection of the control plane. \\\midrule
IaaS security & The provider protects facilities, physical resources, virtualization, core networks, service interfaces, and provider identities while exposing configuration controls to consumers. & Make the shared-responsibility boundary explicit and supply evidence that consumers can use to govern their layers. \\\midrule
Managed platform & The platform standardizes supported languages, runtimes, frameworks, services, deployment units, configuration, scaling, and observability. & Control runtime versions, patching, tenant isolation, build provenance, secrets, service bindings, quotas, and deprecation. \\\midrule
Platform scalability and reliability & The provider scales runtime instances and supporting platform services while hiding much infrastructure management from consumers. & Expose platform limits and dependency behavior; test data, messaging, identity, deployment, and control-plane failure. \\\midrule
PaaS monitoring and security & Platform telemetry and controls must cover application deployment, runtime, dependencies, service identities, data, and consumer actions. & Offer actionable logs and metrics, secure deployment pipelines, isolate tenants, and define which application vulnerabilities remain consumer-owned. \\\midrule
SaaS operating model & The provider manages application code, platform, infrastructure, releases, data services, monitoring, support, and continuity while exposing tenant configuration and use. & Design tenant isolation, release compatibility, configuration governance, data lifecycle, export, identity federation, support access, and incident communication. \\\midrule
SaaS security & Provider control increases, but consumer data, users, entitlements, integrations, and business configuration still determine significant risk. & Supply granular authorization, audit evidence, encryption, secure defaults, retention, deletion, backup, recovery, and transparent subprocessors. \\\midrule
IaaS provisioning decisions & The consumer selects regions, images, instance sizes, networks, storage, access controls, availability patterns, and lifecycle automation. & Use approved templates, policy enforcement, least privilege, patching, inventory, monitoring, backup, cost controls, and deprovisioning. \\\midrule
PaaS provisioning decisions & The consumer chooses runtime and service plans, application topology, scaling, bindings, configuration, identity, data services, and deployment process. & Assess portability, supported versions, quotas, observability, build security, secret handling, data protection, resilience, and platform exit. \\\midrule
SaaS consumption decisions & The consumer maps business processes and data to a provider-operated application and configures roles, policies, integrations, and retention. & Review functional fit, identity, authorization, data location, audit export, availability, support, privacy, customization, lock-in, continuity, and termination. \\\midrule
Provider-consumer responsibility matrix & For every stack layer and lifecycle activity, identify who configures, operates, secures, monitors, backs up, restores, supports, and supplies evidence. & Reconcile the matrix with contracts, architecture, platform behavior, internal staffing, and incident procedures. \\\midrule
\bottomrule
\end{longtable}
\endgroup
\section{Practical Review Checklist}
\begin{longtable}{@{}>{\bfseries\color{Teal}}p{0.16\textwidth}>{\RaggedRight\arraybackslash}p{0.77\textwidth}@{}}
\toprule\textbf{Area} & \textbf{Verification action}\\\midrule\endhead
Scope & $\square$ Define the in-scope services, workloads, users, data, environments, regions, and dependencies for Cloud Delivery Model Considerations. \\\addlinespace[2pt]
Roles & $\square$ Name the provider, consumer, owner, administrator, broker, support, security, finance, and risk responsibilities that apply. \\\addlinespace[2pt]
Boundaries & $\square$ Map organizational, trust, network, tenancy, management, data-location, and failure-domain boundaries. \\\addlinespace[2pt]
Architecture & $\square$ Document how Data centers and resource pools, Scalability and reliability, and supporting mechanisms interact through normal, peak, maintenance, and failed states. \\\addlinespace[2pt]
Identity & $\square$ Verify human and workload identities, privileges, credentials, delegation, administrative access, revocation, and audit evidence. \\\addlinespace[2pt]
Data & $\square$ Trace data collection, transfer, processing, storage, replication, backup, retention, export, and deletion. \\\addlinespace[2pt]
Quality & $\square$ Set measurable capacity, latency, availability, reliability, scalability, and recovery criteria relevant to Monitoring and management. \\\addlinespace[2pt]
Security & $\square$ Challenge the trust and failure assumptions surrounding IaaS security with abuse, bypass, isolation, and compromise scenarios. \\\addlinespace[2pt]
Operations & $\square$ Confirm monitoring, alert ownership, change control, patching, drift management, incident response, and tested recovery. \\\addlinespace[2pt]
Economics & $\square$ Tie resource behavior to meters, units, rates, budgets, anomalies, full-lifecycle cost, and accountable owners. \\\addlinespace[2pt]
Contracts & $\square$ Reconcile architecture behavior with service scope, formulas, exclusions, evidence, notification, remedies, audit rights, and exit terms. \\\addlinespace[2pt]
Verification & $\square$ Preserve reproducible evidence that Managed platform and the overall design satisfy their intended outcomes under realistic load and failure. \\\addlinespace[2pt]
\bottomrule
\end{longtable}
\begin{synthesisbox}
Delivery-model choice changes which party builds, operates, secures, monitors, scales, and pays for each layer. Providers must industrialize shared services and prove quality; consumers must configure and govern the portion they retain. Responsibility shifts upward from IaaS to PaaS to SaaS, but consumer accountability for data, identity, use, risk, and contract fit never disappears.
\par\smallskip
\textbf{Cross-chapter application.} This chapter applies the earlier models and architectures to practical service construction and consumption before cost and quality are quantified.
\end{synthesisbox}
\section{Self-Test Questions}
\begin{enumerate}
\item What central problem does Cloud Delivery Model Considerations address?
\item Define Data centers and resource pools and explain why it matters.
\item Compare Scalability and reliability with Monitoring and management. Which assumptions and responsibilities differ?
\item Trace the normal behavior or decision path associated with IaaS security.
\item Which evidence would you request when reviewing Managed platform?
\item How does this chapter connect to the broader cloud architecture and operating model?
\item Scenario: A team claims that Scalability and reliability is implemented but cannot show current configuration, monitoring, or failure-test evidence. What should the reviewer conclude and do next?
\item Scenario: A workload using Monitoring and management crosses a provider-consumer boundary and then fails during peak demand. Which responsibilities and failure conditions should be investigated?
\item Why is a named mechanism, model, or architecture insufficient evidence that its intended quality or security outcome is achieved?
\item How would you verify that a change involving Platform scalability and reliability improved the intended outcome without shifting risk or cost elsewhere?
\end{enumerate}
\clearpage
\section{Answer Key}
\begin{enumerate}
\item Compare provider and consumer responsibilities, investments, controls, and operating concerns for infrastructure, platform, and software service delivery.
\item The provider invests in physical sites and organizes standardized compute, storage, and network resources into managed pools exposed through virtual infrastructure. The review implication is: Plan capacity, regions, failure domains, lifecycle, utilization, reservations, spares, physical security, and environmental resilience.
\item Replication, clustering, load distribution, elastic allocation, and failover enable variable capacity and continuity. In contrast, Usage, health, audit, SLA, resource, and billing systems operate the service and create consumer-facing evidence. Compare control, ownership, state, dependencies, operating evidence, and failure impact rather than names alone.
\item The provider protects facilities, physical resources, virtualization, core networks, service interfaces, and provider identities while exposing configuration controls to consumers. A trace should name initiating actor, input or signal, decision point, participating components, state change, output, monitoring, and failure handling.
\item Useful evidence should demonstrate the review lens: Control runtime versions, patching, tenant isolation, build provenance, secrets, service bindings, quotas, and deprecation. Include architecture, configuration, event, test, and ownership evidence where applicable.
\item This chapter applies the earlier models and architectures to practical service construction and consumption before cost and quality are quantified.
\item The implementation claim is unsubstantiated. Request the evidence needed to establish current state, then test positive, negative, overload, and failure paths. For this mechanism specifically: Publish realistic limits and test saturation, correlated failure, maintenance, recovery, and performance after loss.
\item Investigate the agreed responsibility split, boundary controls, capacity and quota, workload state, dependent services, provider telemetry, consumer telemetry, escalation, recovery, and whether the contract measures the consumer-visible failure.
\item A label expresses design intent, not runtime proof. Effectiveness depends on placement, configuration, inputs, state, control independence, capacity, monitoring, response, maintenance, and tested behavior during failure.
\item Establish a before-and-after measure, preserve comparable test conditions, verify the expected behavior, inspect adjacent dependencies and costs, and exercise rollback or recovery. Apply this specific lens: Expose platform limits and dependency behavior; test data, messaging, identity, deployment, and control-plane failure.
\end{enumerate}
\section{Key-Term Recap}
\begin{longtable}{@{}>{\RaggedRight\arraybackslash}p{0.27\textwidth}>{\RaggedRight\arraybackslash}p{0.66\textwidth}@{}}
\toprule\rowcolor{Navy}\color{white}\textbf{Term} & \color{white}\textbf{Working meaning}\\\midrule\endhead
Data centers and resource pools & The provider invests in physical sites and organizes standardized compute, storage, and network resources into managed pools exposed through virtual infrastructure. \\\addlinespace[2pt]
Scalability and reliability & Replication, clustering, load distribution, elastic allocation, and failover enable variable capacity and continuity. \\\addlinespace[2pt]
Monitoring and management & Usage, health, audit, SLA, resource, and billing systems operate the service and create consumer-facing evidence. \\\addlinespace[2pt]
IaaS security & The provider protects facilities, physical resources, virtualization, core networks, service interfaces, and provider identities while exposing configuration controls to consumers. \\\addlinespace[2pt]
Managed platform & The platform standardizes supported languages, runtimes, frameworks, services, deployment units, configuration, scaling, and observability. \\\addlinespace[2pt]
Platform scalability and reliability & The provider scales runtime instances and supporting platform services while hiding much infrastructure management from consumers. \\\addlinespace[2pt]
PaaS monitoring and security & Platform telemetry and controls must cover application deployment, runtime, dependencies, service identities, data, and consumer actions. \\\addlinespace[2pt]
SaaS operating model & The provider manages application code, platform, infrastructure, releases, data services, monitoring, support, and continuity while exposing tenant configuration and use. \\\addlinespace[2pt]
SaaS security & Provider control increases, but consumer data, users, entitlements, integrations, and business configuration still determine significant risk. \\\addlinespace[2pt]
IaaS provisioning decisions & The consumer selects regions, images, instance sizes, networks, storage, access controls, availability patterns, and lifecycle automation. \\\addlinespace[2pt]
PaaS provisioning decisions & The consumer chooses runtime and service plans, application topology, scaling, bindings, configuration, identity, data services, and deployment process. \\\addlinespace[2pt]
SaaS consumption decisions & The consumer maps business processes and data to a provider-operated application and configures roles, policies, integrations, and retention. \\\addlinespace[2pt]
Provider-consumer responsibility matrix & For every stack layer and lifecycle activity, identify who configures, operates, secures, monitors, backs up, restores, supports, and supplies evidence. \\\addlinespace[2pt]
\bottomrule
\end{longtable}
\begin{takeawaybox}
Choose and govern a delivery model by tracing every technical and operational layer to a named provider or consumer owner with measurable evidence.
\end{takeawaybox}
\end{document}
